%% ============================================================
%% 第一章 绪论
%% ============================================================
\chapter{绪论}
%% 章标题自动排成“第一章 绪论”（三号黑体居中，每章自动另起一页）

本章介绍选题的理由、课题主要解决的问题，采用的手段、方法，简述研究
成果及其意义。正文汉字为小四号宋体，英文、数字为 Times New Roman，
行距 1.25 倍，每个段落首行自动缩进两个汉字。

正文中引用参考文献应在引用处加方括号注明文献号码，作为上角标，如某某人
对此作了研究\cite{ref1}；也可作为文字段落的组成部分，如数学模型见文献
\cite{ref2}。

\section{研究背景与意义}
%% 节标题自动排成“1.1 研究背景与意义”（小三号黑体居中）

这里是节下的正文内容。

\subsection{国内研究现状}
%% 小节标题自动排成“1.1.1 国内研究现状”（四号黑体左起）
%% 标题层级最多到三级（1.1.1），不要再细分

这里是小节下的正文内容。

%% 小节内的小标题：黑体小四单列一行、缩进两个汉字，序号用 1、2、3……自行书写
\minisection{1. 机构动力学等效集中参数模型}

将机构连续参数（质量和刚度连续分布）离散化，用有限多个自由度表示。

%% 段落标题：宋体加黑小四，置于段首；段内序号依次用⑴⑵⑶、A B C、a b c
（1）\paratitle{建模假设}\quad 建立动力学模型过程中作如下假设：
A.~凸轮机构各构件无制造安装误差；B.~凸轮为刚体；C.~考虑凸轮输入轴扭转
弹性，但输入转速为常数。

\subsection{国外研究现状}

这里是小节下的正文内容。

\section{本文主要研究内容}

这里是节下的正文内容。
